\documentclass[11pt]{article}
\usepackage[margin=1in]{geometry}
\usepackage{amsmath,amssymb}
\usepackage{booktabs}
\usepackage{enumitem}
\usepackage{hyperref}
\usepackage{listings}
\usepackage{xcolor}

\lstset{
  basicstyle=\ttfamily\small,
  keywordstyle=\color{blue},
  commentstyle=\color{gray},
  breaklines=true
}

\title{STeP Functional Emulation Model:\\Design Document}
\author{STeP Team}
\date{April 2026}

\begin{document}
\maketitle

\section{Overview}

The functional emulator computes \emph{values} through a STeP dataflow graph without modeling timing, bandwidth, or hardware resource contention.
Its purpose is to verify that a STeP graph produces the correct output tensor by comparing against a PyTorch reference implementation.

\paragraph{Contrast with cycle-accurate simulation.}
The Rust cycle-accurate simulator models per-cycle tile movement through compute units, on-chip buffers, and off-chip memory channels.
The functional emulator abstracts all of this away: each operator is a pure function from input tensors to output tensors, evaluated once in topological order.

\begin{center}
\begin{tabular}{lcc}
\toprule
& \textbf{Functional Emulator} & \textbf{Cycle-Accurate Sim} \\
\midrule
Language        & Python (PyTorch)  & Rust \\
Semantics       & Value-level       & Cycle-level \\
Output          & Dense tensor      & Dense tensor + timing trace \\
Speed           & Milliseconds      & Seconds--minutes \\
Use case        & Graph correctness & Performance modeling \\
\bottomrule
\end{tabular}
\end{center}

\section{Tensor Representation}

Every node produces a tiled tensor of shape
\[
  (\underbrace{s_0, s_1, \ldots, s_{n-1}}_{\text{stream shape}},\; \text{tile\_r},\; \text{tile\_c})
\]
where the last two dimensions are the spatial tile dimensions and the leading dimensions describe the tile grid (the \emph{stream shape}).

\paragraph{Convention.} Source nodes prepend a leading~1 to the stream shape, i.e.\ $s_0 = 1$, matching the STeP IR convention that every tile stream begins with a singleton batch dimension.

\paragraph{Multi-output nodes.} \texttt{Broadcast} and \texttt{FlatPartition} produce a Python \texttt{list} of tensors rather than a single tensor.
Downstream consumers reference a specific output index.

\section{Execution Model}

\begin{enumerate}[leftmargin=*]
  \item Compute a topological ordering of the graph (via NetworkX).
  \item Walk nodes in order; dispatch each to its handler, which reads predecessor values from a dictionary keyed by instance ID and writes its own result back.
  \item After all nodes are evaluated, the designated \texttt{OffChipStore} node's input is \emph{untiled} into a dense 2-D tensor and returned.
\end{enumerate}

\noindent Untiling reverses the tiling: the leading singleton is stripped, tile grid dimensions are permuted so that row-tiles and col-tiles are adjacent, and the result is reshaped into $(R, C)$.

\section{Operator Semantics}

\subsection{Source Nodes}

\paragraph{\texttt{LinearOffChipLoad}.}
The underlying dense tensor is reshaped into a grid of $(R/\text{tile\_r}) \times (C/\text{tile\_c})$ tiles, flattened, and then indexed by a multi-dimensional stride mapping:
\[
  \text{linear\_idx}(\mathbf{i}) = \sum_{d} i_d \cdot \text{stride}_d
\]
This produces a tensor of shape $(\texttt{out\_shape\_tiled},\;\text{tile\_r},\;\text{tile\_c})$, to which a leading~1 is prepended.

\paragraph{\texttt{LinearOffChipLoadRef}.}
Same tiling as above, but the tile tensor is broadcast over the outer stream dimensions of a \emph{reference} input, enabling weight reuse across batches.

\paragraph{\texttt{SelectGen}.}
Returns the raw multi-hot selection tensor (used by \texttt{FlatPartition}/\texttt{FlatReassemble} for MoE routing).

\subsection{Compute Nodes}

\paragraph{\texttt{UnaryMap}$(f, x)$.}
Applies $f$ element-wise to the full tiled tensor.
Supported functions and their semantics:

\begin{center}
\begin{tabular}{ll}
\toprule
\textbf{Function} & \textbf{Semantics} \\
\midrule
\texttt{Square}         & $x^2$ \\
\texttt{Silu}           & $x \cdot \sigma(x)$ \\
\texttt{Exp}            & $e^x$ \\
\texttt{Rsqrt}          & $1/\sqrt{x}$ \\
\texttt{Pow2}           & $2^x$ \\
\texttt{MulImmediate}$(c)$ & $x \cdot c$ \\
\texttt{AddImmediate}$(c)$ & $x + c$ \\
\texttt{SubImmediate}$(c)$ & $x - c$ \\
\texttt{RowWiseSum}     & $\sum_j x_{ij}$ (reduce last dim, keep dim) \\
\bottomrule
\end{tabular}
\end{center}

\paragraph{\texttt{BinaryMap}$(f, a, b)$.}
Applies $f$ element-wise (with broadcasting) to two tiled tensors.
Supported: \texttt{Matmul} ($ab$ or $ab^\top$), \texttt{Add}, \texttt{Mul}, \texttt{Div}.

\paragraph{\texttt{BinaryMapAccum}$(f, a, b, \text{rank})$.}
Applies the map function $f$ tile-wise, then sum-reduces over the last \texttt{rank} stream dimensions (dimension $-3$ relative to tile dims).
Used for tiled matrix multiplication with accumulation:
\[
  C_{ij} = \sum_k f(A_{ik}, B_{kj})
\]

\paragraph{\texttt{Accum}$(f, x, \text{rank})$.}
Reduces the last \texttt{rank} stream dimensions of $x$ using function $f$:

\begin{center}
\begin{tabular}{ll}
\toprule
\textbf{Function} & \textbf{Semantics} \\
\midrule
\texttt{Add}       & Sum-reduce along dim $-3$ \\
\texttt{Mul}       & Product-reduce along dim $-3$ \\
\texttt{RetileRow} & Concatenate tiles vertically: $(\ldots, K, t_r, t_c) \to (\ldots, K \cdot t_r, t_c)$ \\
\texttt{RetileCol} & Concatenate tiles horizontally: $(\ldots, K, t_r, t_c) \to (\ldots, t_r, K \cdot t_c)$ \\
\bottomrule
\end{tabular}
\end{center}

\subsection{Structural Nodes}

\paragraph{\texttt{Broadcast}$(x, n)$.}
Duplicates the tensor $n$ times; returns a list of $n$ identical references.
Inserted by \texttt{infer\_broadcast} when one output feeds multiple consumers.

\paragraph{\texttt{RepeatStatic}$(x, k)$.}
Inserts a new stream dimension of size $k$ at position $-3$ (before tile dims) and expands.
Used to replicate a reduced value (e.g., row-wise statistics) across $k$ column tiles.

\paragraph{\texttt{Bufferize} / \texttt{Streamify}.}
\texttt{Bufferize} is an identity on values (it marks a buffering boundary in the timing model).
\texttt{Streamify} optionally inserts and expands new stream dimensions with given repeat factors, broadcasting buffered values back into a stream.

\paragraph{\texttt{Flatten}$(\text{min\_rank}, \text{max\_rank})$.}
Merges stream dimensions between \texttt{min\_rank} and \texttt{max\_rank} (counting from the right) into a single dimension.

\paragraph{\texttt{Reshape}$(\text{chunk}, \text{rank}, \text{add\_outer})$.}
Splits the stream dimension at position \texttt{rank} into $(\lceil D / \text{chunk} \rceil, \text{chunk})$, with optional zero-padding and an optional outer singleton dimension.

\paragraph{\texttt{RetileStreamify}$(\text{chunk}, \text{split\_row}/\text{split\_col})$.}
Splits the tile's row or column dimension into chunks, merging the chunk count into the preceding stream dimension.
Optional \texttt{filter\_mask} removes zero-padded chunks.

\paragraph{\texttt{Promote} / \texttt{PromoteOuter}.}
Insert a singleton stream dimension at a specified rank (Promote) or at the outermost position (PromoteOuter).

\paragraph{\texttt{ExpandRef} / \texttt{RepeatRef}.}
Broadcast the input's stream dimensions to match a reference tensor's stream shape.

\subsection{Routing Nodes (MoE)}

\paragraph{\texttt{FlatPartition}$(x, \text{ctrl}, n)$.}
Partitions a flat tile stream into $n$ expert streams using a multi-hot selection tensor from \texttt{SelectGen}.
For each expert $i$, selects tiles where $\text{ctrl}[\cdot, i] > 0$.
Returns a list of $n$ tensors.

\paragraph{\texttt{FlatReassemble}$(\{x_i\}, \text{ctrl})$.}
Inverse of partition: iterates over all token positions, reads from expert streams using per-expert pointers guided by the multi-hot mask, and stacks results.
Output shape follows the control stream convention with a leading~1.

\section{Validation}

The validation harness (\texttt{validate\_functional.py}) compares the functional emulator's output against PyTorch reference implementations for all 14 seed kernels across 55 presets.
The comparison uses relative error:
\[
  \text{rel\_err} = \frac{\|y_{\text{sim}} - y_{\text{ref}}\|_\infty}{\|y_{\text{ref}}\|_\infty + 10^{-12}} < 10^{-5}
\]

All 117 kernel$\times$preset validations pass (as of April 2026).

\section{Comparison with PCL-lite Emulator}

An earlier functional emulator exists in \texttt{AccelOptStep/PCL-lite/step/} (3 files, ${\sim}670$ lines).
It was built for an older version of the STeP IR and targets a different use case: serving as the executable semantics for LLM-based code generation (PCL-lite), where the emulator must validate LLM-synthesized operator graphs.
This section contrasts the two designs.

\subsection{Graph Representation}

\begin{center}
\begin{tabular}{lll}
\toprule
& \textbf{Current (STePTiming)} & \textbf{PCL-lite} \\
\midrule
Graph structure  & NetworkX \texttt{MultiDiGraph}; nodes are & No explicit graph object; \\
                 & \texttt{StepOps} instances with typed edges & \texttt{Stream} objects threaded \\
                 &                                             & through \texttt{OpBase.apply()} calls \\
Traversal        & Topological sort over graph & Implicit via Python call order \\
Node identity    & \texttt{instance\_id} (global counter) & UUID-based names \\
Multi-output     & List of tensors (Broadcast, FlatPartition) & Tuple return (Copy), list (Partition) \\
\bottomrule
\end{tabular}
\end{center}

The current emulator operates on a pre-built graph and walks it; PCL-lite builds the computation inline---each operator call immediately produces a new \texttt{Stream} and (if data is present) evaluates it.

\subsection{Tensor Layout}

\begin{center}
\begin{tabular}{lll}
\toprule
& \textbf{Current} & \textbf{PCL-lite} \\
\midrule
Shape convention & $(\text{stream},\; t_r,\; t_c)$  & $(\text{stream reversed})$; tile dims \\
                 & last two dims are tile spatial     & folded into element dtype \\
Dimension order  & Outermost-first (Python standard) & \textbf{Innermost-first} (reversed) \\
Leading singleton & Always prepended by sources & Not used \\
Data storage     & Single \texttt{torch.Tensor} per node & \texttt{list[torch.Tensor]} per stream \\
                 &                                       & (one tensor per dtype field in STuple) \\
\bottomrule
\end{tabular}
\end{center}

The reversed dimension order in PCL-lite means that \texttt{shape[0]} is the innermost (fastest-varying) stream dimension.
This requires frequent \texttt{[::-1]} reversals when interfacing with PyTorch.
The current emulator uses standard outermost-first ordering, avoiding this indirection.

\subsection{Type System}

PCL-lite has a richer type system that separates element-level types from stream-level structure:

\begin{itemize}[leftmargin=*]
  \item \textbf{Element types}: \texttt{Scalar(dtype)}, \texttt{Tile(dtype, shape)}, \texttt{Multihot(dtype, length)}.
  \item \textbf{Buffer}: wraps an element type with a symbolic shape, representing buffered (non-streaming) data.
  \item \textbf{STuple}: heterogeneous tuple of elements, enabling multi-field streams (e.g., key-value pairs).
  \item \textbf{Symbolic dimensions}: uses \texttt{sympy.Symbol} for shapes that depend on runtime routing decisions.
\end{itemize}

The current emulator has no separate type system---tile dimensions are encoded directly in the tensor shape, and the STeP IR's own \texttt{Tile}/\texttt{DynTile}/\texttt{Buffer} datatype classes (defined in \texttt{step\_py/datatype.py}) carry type metadata.
This is simpler but means the emulator relies on the IR's type annotations rather than enforcing its own type invariants.

\subsection{Operator Coverage}

\begin{center}
\small
\begin{tabular}{lcc}
\toprule
\textbf{Category} & \textbf{Current} & \textbf{PCL-lite} \\
\midrule
\multicolumn{3}{l}{\emph{Compute}} \\
\quad Element-wise map      & UnaryMap, BinaryMap      & Map (via Fn) \\
\quad Map + accumulate      & BinaryMapAccum           & --- \\
\quad Accumulate / reduce   & Accum (Add, Mul, RetileRow/Col) & Accum (arbitrary Fn) \\
\quad Scan (prefix)         & ---                      & Scan \\
\midrule
\multicolumn{3}{l}{\emph{Structure}} \\
\quad Broadcast / copy      & Broadcast                & Copy \\
\quad Repeat                & RepeatStatic, RepeatRef  & Repeat, RepeatRef \\
\quad Bufferize / streamify & Bufferize, Streamify     & Bufferize, Streamify \\
\quad Reshape / flatten     & Reshape, Flatten         & Reshape, Flatten \\
\quad Promote               & Promote, PromoteOuter    & Promote \\
\quad RetileStreamify       & RetileStreamify          & --- \\
\quad ExpandRef             & ExpandRef                & --- \\
\midrule
\multicolumn{3}{l}{\emph{Routing (MoE)}} \\
\quad Partition             & FlatPartition            & Partition \\
\quad Reassemble / merge    & FlatReassemble           & Merge \\
\midrule
\multicolumn{3}{l}{\emph{I/O}} \\
\quad Load                  & LinearOffChipLoad, LoadRef & (external) \\
\quad Store                 & OffChipStore + untiling  & (external) \\
\midrule
\multicolumn{3}{l}{\emph{Misc}} \\
\quad Zip (multi-field)     & ---                      & Zip \\
\quad Identity              & ---                      & Identity \\
\bottomrule
\end{tabular}
\end{center}

Key differences:
\begin{itemize}[leftmargin=*]
  \item PCL-lite's \texttt{Map} is fully generic: the user supplies a \texttt{Fn} object with an \texttt{apply()} method, so any element-wise function can be expressed.
        The current emulator hard-codes each supported function in a dispatch chain.
  \item PCL-lite has \texttt{Scan} (prefix accumulation retaining all intermediate states); the current emulator does not.
  \item PCL-lite has \texttt{Zip} for composing multi-field streams via \texttt{STuple}; the current emulator handles multi-input ops directly in \texttt{BinaryMap}/\texttt{BinaryMapAccum}.
  \item The current emulator has \texttt{BinaryMapAccum} (fused map-then-reduce, critical for tiled matmul), \texttt{RetileStreamify}, and \texttt{ExpandRef}, which PCL-lite lacks.
  \item The current emulator handles I/O nodes (load with stride-based tiling, store with untiling); PCL-lite treats data creation as external to the operator graph.
\end{itemize}

\subsection{Compute Functions}

\begin{center}
\begin{tabular}{ll}
\toprule
\textbf{Current} & \textbf{PCL-lite} \\
\midrule
Fixed set of \texttt{map\_fn} classes & User-defined \texttt{Fn} subclasses \\
(Square, Silu, Matmul, Add, \ldots) & with arbitrary \texttt{apply()} \\
Dispatched by \texttt{isinstance} checks & Dispatched by virtual method call \\
No initial state (Accum fns are also fixed) & \texttt{Fn.getInit()} provides initial accumulator state \\
\bottomrule
\end{tabular}
\end{center}

PCL-lite's approach is more extensible---new functions require only a new \texttt{Fn} subclass.
The current emulator requires adding a branch to \texttt{\_apply\_unary} or \texttt{\_apply\_binary}, but benefits from using PyTorch's vectorized ops on full tiled tensors rather than iterating tile-by-tile.

\subsection{Execution Strategy}

\begin{center}
\begin{tabular}{lll}
\toprule
& \textbf{Current} & \textbf{PCL-lite} \\
\midrule
Evaluation & Batch: full tensor ops via PyTorch & Per-tile: iterates over \\
           & broadcasting and indexing          & \texttt{get\_full\_indices(shape)}, \\
           &                                    & PyTorch within each tile \\
Graph walk & Explicit topological sort           & Implicit (eager, inline) \\
Shape analysis & Not separated from execution   & Dual-phase: shape/type analysis \\
               &                                & runs even without data \\
Symbol handling & Concrete shapes; DynDim for     & \texttt{sympy.Symbol} throughout; \\
               & symbolic MoE dims               & substituted at evaluation time \\
\bottomrule
\end{tabular}
\end{center}

The most significant performance difference: PCL-lite's \texttt{Map} and \texttt{Accum} operators iterate over every tile index individually via \texttt{get\_full\_indices()}, calling \texttt{Fn.apply()} on each tile separately and \texttt{torch.stack}-ing the results.
Both emulators use PyTorch for the per-tile math, but the current emulator applies PyTorch operations to the entire tiled tensor at once (e.g., \texttt{x ** 2} on a tensor of shape $(1, M/t_m, K/t_k, t_m, t_k)$), so PyTorch vectorizes over the tile grid as well---which is orders of magnitude faster for large grids.

\subsection{Design Philosophy}

\begin{center}
\begin{tabular}{lll}
\toprule
& \textbf{Current} & \textbf{PCL-lite} \\
\midrule
Goal        & Validate STeP IR graphs & Executable spec for LLM code synthesis \\
Extensibility & Add IR ops as needed   & Fully generic via Fn/STuple \\
Coupling    & Tightly coupled to \texttt{step\_py} IR & Self-contained (no IR dependency) \\
Performance & Vectorized (fast)       & Scalar iteration (correct but slow) \\
Type safety & Relies on IR types      & Own type system with assertions \\
\bottomrule
\end{tabular}
\end{center}

PCL-lite was designed as a \emph{language definition}: a minimal, self-contained specification of what each operator means, suitable for constraining LLM-generated programs.
The current emulator was designed as a \emph{verification tool}: tightly integrated with the STeP Python frontend, optimized for speed, and validated against 117 kernel$\times$preset configurations.

\end{document}
